\chapter{Document Authoring Guide}
\ChapterLabel{DocumentGuide}
\index{LaTeX}
\index{TikZ}
\index{document conventions}
\index{build tools}
\index{AI agent!LaTeX installation prompt}
\index{Appendix!Document Authoring Guide}

\section{How to Use This Appendix}
Use this appendix to maintain authoring conventions, repository locations,
diagram sources, review rules, tool installation, and PDF build instructions.
Update it whenever the document structure or toolchain changes.

\section{Template Provenance}
\index{template provenance}
This template was developed by \textbf{Darian Muresan, Ph.D.}, in conjunction
with the AI agents \textbf{OpenAI Codex} and
\textbf{Anthropic Claude Code}. Questions, corrections, and proposed
improvements may be sent to
\href{mailto:dmuresan@stevens.edu}{dmuresan@stevens.edu},
\href{mailto:darian@dmmd.net}{darian@dmmd.net}, or
\href{mailto:darian@pictomic.com}{darian@pictomic.com}.

Keep the template attribution distinct from the project-specific document
owner and organization. A derived project may replace the project metadata
while retaining this provenance statement for the underlying template.

\section{AI-Assisted Authoring}
\index{AI agent!architecture workflow}
The authoritative instructions for AI-assisted architecture work, including
task routing, limited context loading, workflows, example prompts, completion
checks, and human-review boundaries, are contained in
\ChapterRef{AIAgentGuide}. This appendix retains installation, build, file, and
LaTeX authoring guidance that applies to both human authors and AI agents.

\section{File Naming}

Every project-owned file created for this document shall begin with the project
file namespace, followed by a Pascal-style responsibility name. The namespace
is defined once in \texttt{tgmManual.tex}. The defining macro is
\texttt{FileNamespace}, and this template sets it to
\texttt{\FileNamespace}. Example chapter files include
\texttt{tgmRequirements.tex}, \texttt{tgmTests.tex}, and
\texttt{tgmValidation.tex}. The convention applies to chapter files, diagram
sources, generated project figures, evidence files, and other project-owned
artifacts. A tool-mandated or externally maintained file name may remain an
explicitly documented exception, but shall not establish a second project
naming convention.

Before creating or renaming files for a derived project, an AI agent shall ask
an authorized human for the project file namespace or confirm a namespace that
is already explicitly configured. The agent shall not invent one. After it is
confirmed, replace \verb|\FileNamespace| and apply the namespace consistently
to every new project-owned file.

\subsection{Document-Authoring Requirements}
The following requirements govern the template itself and every project derived
from it. They are verified by repository inspection and by compiling the
document.

\small
\begin{tabularx}{\linewidth}{L{2.5cm}Y L{3.0cm}}
\caption{Document-Authoring Requirements\label{Table:DocumentAuthoringRequirements}}\\
\toprule
\textbf{ID} & \textbf{Requirement} & \textbf{Verification}\\
\midrule
\endfirsthead
\multicolumn{3}{l}{\small\itshape Document-Authoring Requirements (continued)}\\
\toprule
\textbf{ID} & \textbf{Requirement} & \textbf{Verification}\\
\midrule
\endhead
\bottomrule
\endlastfoot
DOC-FILE-001 & Except for an explicitly documented tool-mandated discovery
file, every new project-owned file shall be named
\texttt{\textless nameSpace\textgreater\textless FileName\textgreater.\textless extension\textgreater},
using the human-approved project namespace and a descriptive Pascal-style file
name. & Inspect all newly created project files.\\
DOC-TIKZ-001 & Every TikZ source shall be stored under
\texttt{TikZ/\textless includingFileBaseName\textgreater/} and shall be named
\texttt{\textless nameSpace\textgreater\textless DrawingName\textgreater.tex}.
The directory name shall exactly match the base name of the LaTeX file that
contains its \texttt{input} command. & Compare each TikZ path with its including
LaTeX file and compile the document.\\
DOC-AI-001 & Before creating or renaming any project-owned file, an AI agent
shall ask for or confirm the project \texttt{\textless nameSpace\textgreater}
and shall apply the confirmed namespace to all files it creates. & Review the
AI interaction record and inspect created file names.\\
DOC-AI-002 & Every repository-level AI instruction file shall point to
\texttt{tgmAIAgentGuide.tex} and the relevant authoritative LaTeX sources,
shall instruct the agent to read them before architecture-related work, and
shall not duplicate their rules. Tool-mandated discovery file names are
documented exceptions to DOC-FILE-001. & Inspect each repository AI instruction
file and compare it with the authoritative LaTeX guidance.\\
\end{tabularx}
\normalsize

\section{One Responsibility per Chapter}

A chapter answers one major engineering question. When content belongs
to another responsibility, link to its authoritative chapter instead of
copying it. In particular, tests define procedures, verification records
requirement compliance, validation evaluates stakeholder use, and release
decides whether the evidence is sufficient.

\section{Adding a Chapter}

Create the namespaced file, begin it with a chapter title and stable chapter
label, add useful index terms, and include it once from
\texttt{tgmManual.tex}. Add or change artifact mappings only in the canonical
traceability matrix.

\begin{verbatim}
\chapter{Example Chapter}
\ChapterLabel{ExampleChapter}
\index{example}
\index{Chapter!Example Chapter}

Chapter content goes here.
\end{verbatim}

\section{Diagram Authoring with TikZ}
\label{Section:TikZAuthoring}
\index{diagram!TikZ source}

TikZ is the diagram-authoring tool for this template. Authors shall use TikZ
for architecture diagrams because the editable source
remains in LaTeX, text uses the document fonts, output stays vector based, and
diagram changes can be reviewed with the surrounding specification. The main
document loads TikZ and a common set of libraries for arrows, positioning,
geometric shapes, fitting, and coordinate calculations.

Every TikZ source shall use the following directory hierarchy:
\begin{center}
\texttt{TikZ/\textless includingFileBaseName\textgreater/\textless nameSpace\textgreater\textless DrawingName\textgreater.tex}
\end{center}
The directory name shall exactly match the base name of the LaTeX file that
contains the diagram's \verb|\input| command. The TikZ file name shall begin
with the configured project namespace and describe the drawing. For example,
the Introduction chapter includes its V-model from:
\begin{center}
\path|TikZ/tgmIntroduction/tgmVModel.tex|
\end{center}
Create additional including-file directories as required; do not place TikZ
sources beside chapter files or in generated-image directories.

The TikZ source file should contain the \texttt{tikzpicture} environment. The
owning chapter shall include that source in a numbered figure, provide a
descriptive caption, place a stable label after the caption, and reference the
figure from the chapter text:

\begin{verbatim}
As shown in Figure~\ref{Figure:HighLevelUseCase}, ...

\begin{figure}[htbp]
  \centering
  \begin{tikzpicture}[
  actor/.style={
    draw=dsnGray,
    fill=dsnGray!8,
    rounded corners=2pt,
    minimum width=2.6cm,
    minimum height=0.8cm,
    align=center,
    font=\small
  },
  use case/.style={
    ellipse,
    draw=dsnBlue,
    fill=dsnLightBlue,
    minimum width=3.8cm,
    minimum height=1.55cm,
    inner sep=1pt,
    align=center,
    font=\small
  },
  association/.style={
    draw=dsnGray,
    line width=0.7pt
  },
  system boundary/.style={
    draw=dsnBlue,
    rounded corners=4pt,
    inner sep=0.45cm
  }
]
  \node[use case] (submit) at (0,0) {\UseCaseRef{UC-001}\\Submit a Service\\Request};
  \node[use case] (status) at (4.25,0) {\UseCaseRef{UC-002}\\View Request\\Status};
  \node[use case] (export) at (8.5,0) {\UseCaseRef{UC-003}\\Export Request\\History};

  \node[system boundary, fit=(submit)(status)(export),
    label={[font=\small\bfseries, text=dsnBlue]above:\ProjectName\ System}]
    (system) {};

  \node[actor, above=1.35cm of submit] (customer) {Customer};
  \node[actor, above=1.35cm of export] (analyst) {Operations Analyst};

  \draw[association] (customer) -- (submit);
  \draw[association] (customer) -- (status);
  \draw[association] (analyst) -- (export);
\end{tikzpicture}

  \caption{High-level use-case diagram.}
  \label{Figure:HighLevelUseCase}
\end{figure}
\end{verbatim}

\section{Building the PDF}

TikZ diagrams compile directly when their source files are included by the
owning chapters. The toolchain requires:

\begin{itemize}
  \item a LaTeX distribution containing TikZ/PGF, \texttt{pdflatex},
  \texttt{bibtex}, and \texttt{latexmk};
  \item the \texttt{makeglossaries} and \texttt{makeindex} commands.
\end{itemize}

\subsection{Installing the Tools on macOS}

The full MacTeX distribution is recommended because it includes the LaTeX tools and packages used by this template. Complete these steps from an administrator account:

\begin{enumerate}
  \item Open Terminal and install Apple's command-line tools, which provide \texttt{make}:
\begin{verbatim}
xcode-select --install
\end{verbatim}
  Complete the installer dialog before continuing.

  \item Install Homebrew by following the current instructions at \url{https://brew.sh/}. Then install MacTeX, which includes TikZ:
\begin{verbatim}
brew install --cask mactex
\end{verbatim}
  MacTeX may instead be downloaded directly from \url{https://tug.org/mactex/}.

  \item Close and reopen Terminal so the MacTeX command-line tools are added to the system path.

  \item Verify that every required command is available:
\begin{verbatim}
command -v make pdflatex bibtex \
  makeglossaries makeindex latexmk
kpsewhich tikz.sty
\end{verbatim}
  Each command shall print a path. A missing line identifies the tool or TikZ
  package that still needs to be installed or added to the system path.
\end{enumerate}

\subsection{Installing the Tools on Windows}

Use a current Windows 10 or Windows 11 computer and complete these steps:

\begin{enumerate}
  \item Download and run the Basic MiKTeX Installer from \url{https://miktex.org/howto/install-miktex}. Choose a private per-user installation when appropriate and select \emph{Always install missing packages on-the-fly} for a computer with Internet access.

  \item Open MiKTeX Console, check for updates, and install them. If \texttt{latexmk} is not found during verification, use the Packages page in MiKTeX Console to install the \texttt{latexmk} package.

  \item Close and reopen PowerShell, then verify the required commands:
\begin{verbatim}
Get-Command pdflatex,bibtex
Get-Command makeglossaries,makeindex,latexmk
kpsewhich tikz.sty
\end{verbatim}
  Each tool shall produce a command record. A command-not-found error identifies a missing installation or system-path entry.
\end{enumerate}

On an organization-managed computer, authors shall ask the system administrator
before installing software or changing a managed execution policy.

\subsection{Installing the Tools on Ubuntu or Debian Linux}

Open a terminal and install TikZ/PGF, GNU Make, Latexmk, and a TeX
Live set that contains the other packages used by this template:

\begin{verbatim}
sudo apt update
sudo apt install make latexmk texlive-pictures \
  texlive-latex-extra \
  texlive-fonts-recommended texlive-bibtex-extra \
  texlive-extra-utils
\end{verbatim}

Verify the required commands:

\begin{verbatim}
command -v make pdflatex bibtex \
  makeglossaries makeindex latexmk
kpsewhich tikz.sty
\end{verbatim}

For another Linux distribution, install the equivalent TikZ/PGF, GNU
Make, Latexmk, and full or extended TeX Live packages with its package manager.

\subsection{AI-Agent Installation}
\label{Section:AIAgentInstallation}

An AI coding agent can inspect the host, identify missing commands, request
approval for system changes, install the appropriate packages, and verify the
build. Give the agent the project directory and the following reusable request:

\begin{verbatim}
In this LaTeX project, inspect the build files and determine
which operating system and package manager are available.
Verify TikZ with `kpsewhich tikz.sty`, and verify make,
pdflatex, bibtex, makeglossaries, makeindex, and latexmk.
Install only missing tools using the
platform's standard package manager. Ask for approval before
administrator-level, system-wide, or networked installation.
Do not replace project sources or remove user files. After
installation, compile included TikZ sources, run the full platform
build, inspect the log for errors, unresolved
references, and overfull boxes, and report the generated PDF.
\end{verbatim}

If the agent cannot install software, ask it to report the exact missing
command, the attempted verification, and platform-specific installation
instructions. Never give an agent an unrestricted instruction to delete build
directories or modify system security policy.

\subsection{Building on macOS and Linux}

From the template directory, run:

\begin{verbatim}
make -f makefileMacLinux pdf
\end{verbatim}

Use \texttt{make -f makefileMacLinux clean} to remove generated build files while preserving authored sources and the final PDF.

\subsection{Building with Windows PowerShell}

From PowerShell in the template directory, run:

\begin{verbatim}
.\makefilePowerShell.ps1 -Action pdf
\end{verbatim}

Use \texttt{-Action clean} in place of \texttt{-Action pdf} to remove generated build files. If local execution policy prevents scripts from running, use an approved organizational policy or allow scripts only for the current PowerShell process:

\begin{verbatim}
Set-ExecutionPolicy -Scope Process Bypass
.\makefilePowerShell.ps1 -Action pdf
\end{verbatim}

\subsection{Manual Build Sequence}

When the platform build script is unavailable, run the following commands from
the template directory. The final two LaTeX passes resolve tables of contents,
citations, glossary links, the index, and cross-references.

Included TikZ sources require no separate rendering command; \texttt{pdflatex}
compiles them with their owning chapters:

\begin{verbatim}
pdflatex -interaction=nonstopmode -halt-on-error tgmManual.tex
bibtex tgmManual
makeglossaries tgmManual
makeindex tgmManual.idx
pdflatex -interaction=nonstopmode -halt-on-error tgmManual.tex
pdflatex -interaction=nonstopmode -halt-on-error tgmManual.tex
\end{verbatim}

\subsection{Build Verification and Troubleshooting}

Confirm that the output PDF opens, all expected chapters and appendices appear,
the glossary and index are populated, hyperlinks reach their targets, and no
table or figure is clipped. Search \texttt{tgmManual.log} for fatal errors,
undefined references, missing files, and overfull boxes. Common remedies are:

\begin{itemize}
\item open a new terminal after installation so its search path is refreshed;
\item on macOS, verify that \texttt{/Library/TeX/texbin} is on the search path;
\item rerun the complete build after adding citations, glossary entries, index
entries, figures, labels, or references;
\item compile and inspect the owning chapter after changing a TikZ source;
\item install the named LaTeX package when the log reports a missing
\texttt{.sty} file.
\end{itemize}

\section{Authoring Rules}

\begin{itemize}
\item Add the newest approved change to the top of the document history.
\item Introduce specialized terms through the glossary and use them consistently.
\item Give every controlled artifact a stable explicit ID.
\item Use \texttt{shall} only for binding requirements.
\item Keep normative requirements separate from explanation and examples.
\item Apply the approved project namespace to every new project-owned file as
specified in \hyperref[Table:DocumentAuthoringRequirements]{the
document-authoring requirements}.
\item Use TikZ for all diagrams. Store every source in
\path|TikZ/<includingFileBaseName>/| and name it
\path|<nameSpace><DrawingName>.tex|.
\item Include each TikZ source from its owning chapter in a numbered figure
with a caption and stable label, and reference that figure in the chapter text.
\item Compile and inspect the PDF after structural or table changes.
\end{itemize}

\section{Citations and Bibliography}

Store controlled reference records in \texttt{tgmReferences.bib}. Cite a
source from a chapter with a stable BibTeX key:

\begin{verbatim}
\cite{ISO29148}
\end{verbatim}

The main file applies the style and prints one document-level bibliography:

\begin{verbatim}
\bibliographystyle{IEEEtran}
\bibliography{tgmReferences}
\end{verbatim}

Both platform build scripts run BibTeX automatically.

Bibliography entries support claims and definitions; they do not replace
project-owned requirements, architecture decisions, test procedures, or
approval evidence.

\section{Selective Builds}

For a large project, LaTeX's \verb|\includeonly| command may select chapters
during drafting. The main file contains a commented example:

\begin{verbatim}
% \includeonly{
%   prologue,
%   tgmRequirements,
%   tgmTests
% }
\end{verbatim}

Remove the leading comment characters and list include names without the
\texttt{.tex} extension. LaTeX can retain cross-references and citations from
omitted chapters by reading their existing auxiliary files, but those files
may be absent or stale. Selective builds are therefore drafting aids only.
Hyperref may also report missing destinations for links into deliberately
omitted chapters; these draft warnings shall disappear in the full build.
Comment out the complete \verb|\includeonly| block and run a full build before
verification or release so missing chapters, references, glossary entries,
citations, and evidence cannot be hidden.
